% Part: first-order-logic
% Chapter: syntax-and-semantics
% Section: Structures

\documentclass[../../../include/open-logic-section]{subfiles}

\begin{document}

\olfileid[hi]{fol}{syn}{str}

\olsection{प्रथम-क्रम भाषाओं के लिए \printtoken{P}{structure}}

\begin{explain}
प्रथम-क्रम भाषाएँ अपने-आप में \emph{अनिर्वचित} होती हैं: !!{constant}s,
!!{function}s और !!{predicate}s के साथ कोई विशिष्ट अर्थ जुड़ा नहीं होता।
\article{structure} \emph{!!{structure}} निर्दिष्ट करके उन्हें अर्थ दिया
जाता है। संरचना \emph{प्रांत} निर्दिष्ट करती है, अर्थात् वे वस्तुएँ जिन्हें
!!{constant} चुनते हैं, जिन पर !!{function} कार्य करते हैं और जिन पर
परिमाणक विचरण करते हैं। इसके अतिरिक्त, वह बताती है कि कौन-से !!{constant}
कौन-सी वस्तुएँ चुनते हैं, कोई !!a{function} वस्तुओं को वस्तुओं पर किस प्रकार
प्रतिचित्रित करता है, और !!{predicate} किन वस्तुओं पर लागू होते हैं।
!!^{structure}s तर्कशास्त्र की \emph{अर्थगत} धारणाओं---जैसे तात्पर्य, वैधता और
संतुष्टनीयता---का आधार हैं। साहित्य में इन्हें अलग-अलग जगह ``संरचनाएँ'',
``निर्वचन'' या ``प्रतिरूप'' कहा जाता है।
\end{explain}

\begin{defn}[!!^{structure}s]
\Article{structure} \emph{!!{structure}}~$\Struct M$, प्रथम-क्रम तर्क की
भाषा~$\Lang{L}$ के लिए, निम्न अवयवों से बनी होती है:
\begin{enumerate}
\item \emph{प्रांत:} एक अरिक्त समुच्चय, $\Domain M$
\item \emph{!!{constant}s का निर्वचन:} प्रत्येक !!{constant}~$c$ के लिए,
  जो~$\Lang{L}$ में है, !!a{element} $\Assign{c}{M} \in \Domain M$
\item \emph{!!{predicate}s का निर्वचन:} प्रत्येक $n$-स्थानी
  !!{predicate}~$R$ के लिए, जो~$\Lang{L}$ में है ($\eq$ को छोड़कर), एक
  $n$-स्थानी संबंध
  $\Assign{R}{M} \subseteq \Domain{M}^n$
\item \emph{!!{function}s का निर्वचन:} प्रत्येक $n$-स्थानी
  !!{function}~$f$ के लिए, जो~$\Lang{L}$ में है, एक $n$-स्थानी फलन
  $\Assign{f}{M}
  \colon \Domain{M}^n \to \Domain{M}$
\end{enumerate}
\end{defn}

\begin{ex}
अंकगणित की भाषा के लिए !!^a{structure}~$\Struct M$ में एक समुच्चय, उस
समुच्चय $\Domain M$ का एक अवयव $\Assign{\Obj 0}{M}$ जो
!!{constant}~$\Obj 0$ का निर्वचन है, एक एक-स्थानी फलन
$\Assign{\Obj \prime}{M} \colon \Domain{M} \to \Domain M$, दो दो-स्थानी
फलन $\Assign{\Obj +}{M}$ और $\Assign{\Obj
    \times}{M}$---दोनों
$\Domain M^2 \to \Domain M$---और एक दो-स्थानी संबंध
$\Assign{\Obj <}{M} \subseteq \Domain{M}^2$ होते हैं।

ऐसी संरचना का एक सहज उदाहरण निम्न है:
\begin{enumerate}
\item $\Domain N = \Nat$
\item $\Assign{\Obj 0}{N} = 0$
\item $\Assign{\Obj \prime}{N}(n) = n + 1$, सभी $n \in \Nat$ के लिए
\item $\Assign{\Obj +}{N}(n, m) = n + m$, सभी $n, m \in \Nat$ के लिए
\item $\Assign{\Obj \times}{N}(n, m) = n\cdot m$, सभी $n, m \in \Nat$ के लिए
\item $\Assign{\Obj <}{N} = \Setabs{\tuple{n, m}}{n \in \Nat, m \in
  \Nat, n < m}$
\end{enumerate}
इस प्रकार परिभाषित संरचना~$\Struct N$ को,~$\Lang L_A$ के लिए,
\emph{अंकगणित का मानक प्रतिरूप} कहते हैं, क्योंकि वह~$\Lang L_A$ के
अतार्किक अचरों का निर्वचन ठीक अपेक्षित ढंग से करती है।

फिर भी~$\Lang
L_A$ के लिए अनेक अन्य !!{structure}s संभव हैं। उदाहरणतः पूर्णांकों का
समुच्चय~$\Int$ प्रांत के रूप में,~$\Nat$ के स्थान पर लेकर हम $\Obj
0$, $\Obj \prime$, $\Obj +$, $\Obj \times$, $\Obj <$ के निर्वचन उसी
अनुसार परिभाषित कर सकते हैं। हम~$\Lang L_A$ के लिए ऐसी संरचनाएँ भी
परिभाषित कर सकते हैं जिनका संख्याओं से दूर-दूर तक कोई संबंध न हो।
\end{ex}

\begin{ex}
किसी संरचना~$\Struct M$ को, समुच्चय सिद्धांत की भाषा~$\Lang L_Z$ के लिए,
केवल एक समुच्चय और एक दो-स्थानी संबंध चाहिए। अतः तकनीकी रूप से, उदाहरण के
लिए, व्यक्तियों के समुच्चय के साथ ``$x$, $y$ से अधिक आयु का है'' संबंध को
$\Lang L_Z$ की !!a{structure} बनाया जा सकता है; इसी प्रकार~$\Nat$ के साथ
$n \ge m$ को भी, जहाँ $n, m \in \Nat$।

$\Lang L_Z$ की एक विशेष रोचक !!{structure} वह है जिसमें प्रांत के
!!{element}s वास्तव में समुच्चय हैं और $\Obj \in$ का निर्वचन सचमुच ``$x$,
~$y$ का !!a{element} है'' संबंध है। यह \emph{आनुवंशिक रूप से परिमित
समुच्चयों} की !!{structure}~$\Struct{{HF}}$ है:
\begin{enumerate}
\item $\Domain{{{HF}}} = \emptyset \cup \Pow{\emptyset} \cup
  \Pow{\Pow{\emptyset}} \cup \Pow{\Pow{\Pow{\emptyset}}} \cup \dots$;
\item $\Assign{\Obj \in}{{{HF}}} = \Setabs{\tuple{x, y}}{x, y \in
  \Domain{{{HF}}}, x \in y}$.
\end{enumerate}
\end{ex}

\begin{digress}
किसे !!a{structure} मानना है, इस बारे में रखी गई शर्तें हमारे तर्क को
प्रभावित करती हैं। उदाहरणतः, रिक्त प्रांतों को निषिद्ध करने से परिमाणित
वाक्यों की संतुष्टि (या सत्यता) के सामान्य विवरण के अधीन यह सुनिश्चित होता
है कि $\lexists[x][(!A(x) \lor \lnot !A(x))]$ वैध है---अर्थात् तार्किक सत्य
है। और यह शर्त कि सभी !!{constant}s को प्रांत की किसी वस्तु का निर्देश करना
ही चाहिए, सुनिश्चित करती है कि अस्तित्वपरक सामान्यीकरण एक उचित अनुमान-रूप
है: $!A(a)$, अतः $\lexists[x][!A(x)]$। यदि हम नामों को प्रांत के बाहर की
वस्तुओं का निर्देश करने, अथवा किसी का भी निर्देश न करने, देते तो हम
\emph{मुक्त तर्क} की ओर बढ़ते। उसमें अस्तित्वपरक सामान्यीकरण के लिए एक
अतिरिक्त आधारवाक्य चाहिए: $!A(a)$ और $\lexists[x][\eq[x][a]]$, अतः
$\lexists[x][!A(x)]$।
\end{digress}

\end{document}
